\documentclass[a4paper, 12pt]{article}

\usepackage[a4paper, top=12mm, bottom=12mm, left=8mm, right=12mm]{geometry}

\usepackage{polyglossia}
\setdefaultlanguage[babelshorthands=true]{russian}
\setotherlanguage{english}

\usepackage{fontspec}
\setmainfont{FreeSerif}
\newfontfamily{\russianfonttt}[Scale=0.7]{DejaVuSansMono}

\usepackage[tiny, compact]{titlesec}

\usepackage{titling}
\setlength{\droptitle}{-1cm}
\setlength{\parskip}{1.5ex}
\pretitle{\begin{center}\begin{bfseries}\Large}
\posttitle{\par\end{bfseries}\end{center}}
\preauthor{\begin{center}\normalsize}
\postauthor{\par\end{center}\vspace{-1.8cm}}

\usepackage{enumitem}
\setlist[enumerate]{topsep=0pt, itemsep=0.2pt, leftmargin=2em}
\setlist[itemize]{topsep=0pt, itemsep=0.2pt, leftmargin=2em}

\usepackage{hyperref}
\usepackage{bookmark}
\usepackage{csquotes}
\usepackage{amsmath}
\usepackage{graphicx}
\usepackage{float}
\usepackage{booktabs}


\title{Реализация свёртки изображений}

\author{Меньшикова Елизавета Фёдоровна}

\date{}

\begin{document}

\maketitle

\begin{flushright}
    Группа: \emph{25.Б42-мм}

    Кафедра: \emph{системного программирования СПбГУ}

    Научный руководитель: \emph{Пономарев Николай Алексеевич}
    
    Номер семестра практики: \emph{2}
\end{flushright} 

\section{Введение}

Свёртка --- это операция над двумя функциями, широко применяемая, в частности, в цифровой обработке сигналов и обработке изображений\footnote{\tiny LAMAGraph. Image Convolution. URL: \url{https://lamagraph.github.io/intro-to-fpga-with-clash/Appendix-ImageConvolution/index.html} (дата обращения: 21.08.2026).}. Для дискретных функций одномерная свёртка определяется как:

\begin{equation*}
    (f * g)[n] = \sum_{m=-\infty}^{\infty} f[m] \cdot g[n - m]
\end{equation*}

Близкой к свёртке операцией является кросс-корреляция. Обе операции применяют ядро к локальным областям изображения, но при свёртке используется ядро, отражённое относительно центра, а при кросс-корреляции --- исходное. Для простоты восприятия в данной работе используется кросс-корреляция. Такой подход также применяется в практических библиотеках обработки изображений: например, функция filter2D в OpenCV вычисляет корреляцию\footnote{\tiny OpenCV Documentation. imgproc::filter2D. URL: \url{https://docs.opencv.org/4.5.5/d4/d86/group__imgproc__filter.html} (дата обращения: 21.08.2026).}.

Для дискретного изображения $I$ и ядра $K$ размера $k \times k$ выходное значение пикселя вычисляется по формуле:

\begin{equation*}
    O[i, j] = \sum_{u=0}^{k-1} \sum_{v=0}^{k-1} I_{\text{pad}}[i + u, j + v] \cdot K[u, v],
\end{equation*}

где:
$O[i,j]$~--- интенсивность пикселя выходного изображения; $I_{\text{pad}}$~--- исходное изображение с искусственно добавленной граничной областью (padding).

При обработке граничных пикселей часть области, необходимой для вычисления, выходит за пределы изображения, поэтому необходимо определить способ обработки границ.

В рамках учебной практики реализация позволяет изучить работу с изображениями и многомерными массивами, реализацию математического алгоритма, обработку граничных случаев, автоматическое тестирование и сравнение производительности различных реализаций.

\section{Постановка задачи}

\textbf{Целью работы является} реализация алгоритма двумерной свёртки изображений и сравнительный анализ её производительности с библиотечными решениями (OpenCV, Pillo).

\textbf{Для достижения цели требуется решить следующие задачи:}

\begin{enumerate}
    \item Выбрать подходящие инструменты и библиотеки для реализации проекта.
    \item Реализовать алгоритм свёртки с поддержкой различных ядер и режимов обработки граничных пикселей.
    \item Разработать интерфейс командной строки для управления параметрами программы.
    \item Провести тестирование корректности работы программы по эталонным значениям.
    \item Выполнить эксперименты по сравнению производительности собственной реализации с библиотечными решениями.
\end{enumerate}

\section{Реализация}

Для разработки выбран Python 3.11, так как он предоставляет необходимые библиотеки для работы с массивами (NumPy) и изображениями (Pillow). Управление зависимостями настроено через pyproject.toml (сборщик hatchling). Все настройки проекта хранятся в одном файле: инструменты проверки качества (Ruff, Mypy, Pytest) читают параметры напрямую из него, что упрощает настройку CI и обеспечивает одинаковую сборку на любом компьютере.

Алгоритм свёртки реализован через вложенные циклы. Обработка границ выполнена через векторизованные операции над индексами NumPy: формируются полные массивы координат, которые трансформируются за одну операцию без перебора точек в цикле.

Реализованы 8 типов ядер и 4 режима обработки краёв. Выбранные ядра и способы обработки границ доступны в авторской реализации, OpenCV и Pillow.

\subsection{Запуск программы}

Выбор входного и выходного файла, типа ядра и режима обработки краёв осуществляется через интерфейс командной строки. Интерфейс командной строки сделан на Typer\footnote{ Typer URL:\url{https://typer.tiangolo.com/} (дата обращения: 21.08.2026).}. Изначально использовался argparse, но было решено перейти на Typer, как на более современную версию, позволяющую сократить и упростить код.

Пример запуска программы:
\begin{verbatim}
python -m image_convolution.main input.png output.png --kernel box_blur --border constant
\end{verbatim}

Где:
\begin{itemize}
    \item \texttt{input.png}~--- путь к входному изображению;
    \item \texttt{output.png}~--- путь для сохранения результата;
    \item \texttt{--kernel}~--- название ядра;
    \item \texttt{--border}~--- тип границы.
\end{itemize}

\subsection{Тестирование}

Корректность  работы алгоритма проверяется интеграционными тестами (pytest). В качестве эталонов используются изображения, полученные ранее той же функцией свёртки. Сравнение производится с абсолютным допуском = 1.0, что допускает незначительные различия из-за особенностей вычислений с плавающей точкой.

\section{Эксперименты}

Измерения проводились на системе со следующими характеристиками:

\begin{itemize}

\item \textbf{Процессор:} AMD Ryzen 7 8745HS w/ Radeon 780M Graphics (архитектура x86\_64).

\item \textbf{Оперативная память:} 12 ГБ.

\item \textbf{Операционная система:} Ubuntu 24.04.4 LTS.

\end{itemize}

Эксперименты выполнялись в стандартном окружении операционной системы с зафиксированной тактовой частотой процессора (режим performance). Для снижения фонового шума были отключены пользовательские приложения. Данный подход обеспечивает воспроизводимость результатов в типовом пользовательском окружении, однако допускает большую вариативность измерений по сравнению со специализированными лабораторными условиями, поскольку не применялась аппаратная изоляция ядер CPU.

\textbf{Входные данные: } Для всех трёх реализаций использовались идентичные входные массивы, созданные функцией numpy.random.rand(size, size, 3).astype(np.float32). Это гарантирует, что каждая функция получает одинаковую матрицу случайных значений типа float32 в диапазоне [0,1], что исключает влияние различий в генерации данных на результаты сравнения. Использовались квадратные изображения размеров $128\times 128$, $512\times 512$, $1024\times 1024$ и $2048\times 2048$ пикселей. Применение синтетического равномерного шума минимизирует влияние пространственной локальности пикселей, характерной для реальных фотографий.

\textbf{Параметры свёртки: } ядро box\_blur ($3\times 3$), граница constant.

\textbf{Порядок измерений:} Тестирование проводилось с использованием pytest-benchmark с предварительным прогревом. Результаты измерений сохранялись в JSON-файл и далее обрабатывались отдельно. Измерения проводились для авторской реализации, OpenCV и Pillow в цветном и чёрно-белом режимах.

Для обработки результатов вычислялись среднее время выполнения и стандартное отклонение. Также рассчитывалось отношение времени авторской реализации к OpenCV и Pillow.

\begin{table}[h]

\centering

\caption{Сравнение производительности реализаций свёртки}

\label{tab:performance}

\begin{tabular}{llrrr}

    \toprule

    \textbf{Режим} & \textbf{Размер} & \textbf{Авторская (мс)} & \textbf{OpenCV (мс)} & \textbf{Pillow (мс)} \\

    \midrule

    Цветной & 1024 px & $7527{,}4 \pm 20{,}6$ & $0{,}636 \pm 0{,}075$ & $16{,}49 \pm 0{,}05$ \\

    Цветной & 2048 px & $31184{,}3 \pm 88{,}3$ & $26{,}12 \pm 0{,}93$ & $91{,}05 \pm 0{,}14$ \\

    Цветной & 4096 px & $123938{,}8 \pm 580{,}5$ & $61{,}45 \pm 1{,}93$ & $367{,}56 \pm 0{,}76$ \\

    Полутоновой & 1024 px & $2614{,}1 \pm 10{,}3$ & $0{,}088 \pm 0{,}027$ & $4{,}87 \pm 0{,}07$ \\

    Полутоновой & 2048 px & $10550{,}8 \pm 63{,}7$ & $1{,}120 \pm 0{,}082$ & $24{,}34 \pm 0{,}13$ \\

    Полутоновой & 4096 px & $41971{,}3 \pm 131{,}6$ & $27{,}04 \pm 0{,}84$ & $107{,}61 \pm 0{,}61$ \\

    \bottomrule

\end{tabular}

\end{table}

Авторская реализация медленнее OpenCV примерно в 1194 --- 29800 раз и Pillow в 337 --- 537 раз в зависимости от размера и режима изображения.

При увеличении стороны изображения в 4 раза количество пикселей увеличивается в 16 раз. Время работы авторской реализации увеличивается примерно в 16 раз, что соответствует теоретической сложности алгоритма $O(height\times width)$ при фиксированном размере ядра.

Полученные различия во времени выполнения авторской реализации и библиотечных решений являются статистически значимыми для всех рассмотренных размеров и режимов ($p < 0{,}05$).

Значительное отставание авторской реализации связано с тем, что вычисления выполняются на Python, тогда как OpenCV и Pillow используют оптимизированные низкоуровневые реализации. Основные операции обработки изображений в них выполняются в скомпилированном коде, что существенно уменьшает накладные расходы.

\section{Заключение}

В ходе выполнения работы были достигнуты следующие цели:

\begin{enumerate}
    \item Изучен алгоритм двумерной свёртки и методы обработки границ изображений.
    \item Реализована собственная версия свёртки с поддержкой различных ядер и режимов padding на Python.
    \item Реализованы CLI-интерфейс и набор тестов для собственной реализации.
    \item  Проведён сравнительный анализ производительности собственной реализации с библиотечными решениями.
\end{enumerate}

Репозиторий с реализацией: \url{https://github.com/MenshikovaElizaveta/img-conv}

\end{document}
